\section{Workflow: does it fit how I actually work?}
\label{sec:workflow}

The best safe tool fails if it does not fit. The eighth question asks whether the
system reduces burden or adds to it, and how it changes consent and the clinic day,
because a tool that fights the workflow will be worked around or abandoned. This is
not a soft concern. Clinicians discard even useful guidance once it has disappointed
them, and they do so more readily than they discard a fallible human colleague; the
phenomenon of algorithm aversion shows that confidence, once lost to a visible
error or to daily friction, is hard to recover \cite{dietvorst2015algorithm}. A
system that adds clicks, interrupts the consult, or duplicates documentation will
lose the clinician's trust regardless of verifications.

The answer is workflow integration: the system is embedded into the existing steps
of the clinic day rather than bolted alongside them. Each action still passes verification before generation and the ten-gate VVUQ examination, but the
clinician experiences this as a recommendation that arrives, with its
reasoning and an ACCEPT, ESCALATE, or BLOCK already attached, and a countersignature
that doubles as the documentation. Consent is handled the same way, as a verified
process woven into the encounter rather than a separate hurdle; the role of an
honest, well-supported consent conversation in earning patient trust is well
established \cite{novak2020patient}, and the author's patient-journey lineage traces
a single patient moving from referral through consent to follow-up without the
workflow ever leaving the clinician's normal path \cite{Kawchak2026CancerPatientsJourney}.

\autoref{tab:workflow} sets out the clinic touchpoints and contrasts each one
without the system and with it, so the clinician can see where burden is removed
rather than added. Referral and eligibility, informed consent, the clinic decision,
documentation, and follow-up monitoring each become a step where the verified record
is produced as a by-product of the work rather than as extra work. Follow-up in
particular is where continuous monitoring lives: the same drift and uncertainty
controls that protect reliability run quietly after the visit, surfacing only when a
patient needs attention.

Figure 10 places this single clinic encounter in its larger arc, from one bedside
decision to site-level practice to trial-wide calibrated trust, meaning reliance
proportioned to demonstrated capability. The author's national-platform work shows
how that arc scales beyond one clinic while the workflow at each bedside stays
recognizable \cite{Kawchak2026NationalPlatformPhysicalAI}. The point of the scaling
is not uniformity for its own sake but that the same verified,
experience a clinician trusts at one bedside is the experience an entire trial can
be built on.

 

\begin{cfwfig}
\node[cfone] (a) at (0,0) {Bedside}; \node[cftwo] (b) at (4.0,0) {Site};
\node[cfact] (c) at (8.0,0) {Trial}; \node[cfhope] (d) at (12.0,0) {Platform trust};
\draw[cfedge] (a)--(b); \draw[cfedge] (b)--(c); \draw[cfedge] (c)--(d);
\end{cfwfig}
\vspace{-0.8cm}
\cfwcaption{Figure 10. From bedside to trial-wide trust: a single verified bedside decision scales through the site and the
trial.}

\needspace{9\baselineskip}\begin{table}[H]
\footnotesize
\begin{tabularx}{\textwidth}{@{}L{3.3cm} Y L{6.3cm}@{}}
\toprule
\textbf{Workflow touchpoint} & \textbf{Without the system} & \textbf{With the system}\\
\midrule
Referral and eligibility & Manual chart review vs.\ criteria, easy delayed & Eligibility checked, recorded as referral arrives\\
Informed consent & A separate hurdle bolted onto the visit & Verified consent process woven into encounter\\
The clinic decision & Judgment made, documented from memory & Recomm. in place with reasoning and gate result\\
Documentation & Extra notes written after the fact & A countersignature that is itself the audit record\\
Follow-up monitoring & Periodic, manual, easy to miss & Continuous monitoring surfaces when needed\\
\bottomrule
\end{tabularx}
\caption{The clinic day at each workflow touchpoint, contrasted without and with the
system, showing where workflow integration removes burden rather than adding it.}
\label{tab:workflow}
\end{table}

Workflow is answered by fit: a system embedded in the clinic day and the consent
process, measured for burden, and scaled from one bedside decision to trial-wide
calibrated trust. Workflow integration and continuous monitoring are the concrete
provisions that make the system something a clinician uses rather than works around.

The eight questions are not eight separate tests but eight doors into the same room.
Each one opens onto a specific mechanism, a gate, a record, a role, or a control,
and the clinician who has walked through all eight has the evidence to calibrate
trust rather than to grant or withhold it wholesale. \autoref{tab:synthesis}
gathers the eight questions and the mechanism that answers each, so the framework
can be read at a glance and carried into a credentialing committee or a tumor board.

\needspace{12\baselineskip}\begin{table}[H]
\footnotesize
\begin{tabularx}{\textwidth}{@{}L{3.0cm} Y@{}}
\toprule
\textbf{Question} & \textbf{The mechanism that answers it}\\
\midrule
Competence & Local validation and an assurance run, with 172 of 172 gate tests passed before action\\
Safety & The ten-gate VVUQ examination and the catastrophe-predicate BLOCK before the patient\\
Transparency & Explainable output and a complete, hash-chained audit trail\\
Oversight & Human-in-the-loop control, the ESCALATE path, and an unconditional override\\
Equity & Subgroup performance reporting and continuous bias surveillance\\
Reliability & Reproducibility checks, drift monitoring, and quantified uncertainty\\
Accountability & Named roles and a signed standard operating procedure\\
Workflow & Integration into the clinic day and consent, with continuous monitoring\\
\bottomrule
\end{tabularx}
\caption{The eight confidence questions and the specific mechanism that answers each.}
\label{tab:synthesis}
\end{table}

A clinician who can answer all eight questions has what calibrated trust requires:
not faith, but the evidence to rely on the system exactly as far as it has earned,
and no further. The three questions that opened this framework now resolve into one.
Would I let this system near my own family member, could I defend the decision at
tumor board, and could I show an auditor exactly what happened and why? Each can be
answered yes, because a gate, a record, a role, or a control stands behind it. That
is the adoption decision this framework asks the clinician to make, and to sign.
